
\begin{table}[!htbp]\centering\caption{模型口径与检验层次}\label{tab:scope}
\begin{tabularx}{\linewidth}{p{.15\linewidth}XX}\toprule
对象 & 实现与依据 & 结果的统一解释\\\midrule
历史分类 & 年鉴“一次电力及其他能源”保持原分类；前瞻以七类电源构造非化石工程映射 & 历史第四分类与模型非化石结果分别命名，分类映射用于条件比较\\
工程与行为 & 全国聚合母线、连续机组、外生水电可用率；建设、退役、成本及响应部分校准 & 217维为状态维数；12类资产承担投资选择，小时出力按聚合母线与连续爬坡执行\\
进口折标 & 煤、原油和天然气用官方参考系数，成品油以燃料油系数作聚合代表 & 基年满足$rP+I=C$；平衡调节项单列，折标低档在物理域内修订为$-2\%$\\
历史检验 & 固定起点、给定已实现GDP的共享预期与资产组件检验 & 用于组件误差诊断；独立进口映射从2024年前瞻起点生效\\
机制比较 & 主实验保留自主趋势、关闭主动响应；冻结份额另作诊断 & 主动反馈贡献与外生载能趋势分开解释；较高近期需求外推只作需求压力对照\\
尺度与信息 & 年度内部闭环，关键年状态单向下钻；原生试验24小时前瞻，信息挑战只改变事件可见时长 & 滚动控制采用24小时优化窗；离线参照固定实际前48小时后一次优化312小时；年度投资仍由年度规则产生\\
事件与恢复 & 工程天气、负荷和到货；连续继承燃料、SOC、爬坡记忆和债务 & 对应固定设计事件集合；库存以事件前、固定初始及同小时无事件分别参照，格点数描述工程条件覆盖\\
费用与搜索 & 按模拟订单付资；工程资本、小时资源费用、合同现金分账；固定预算搜索 & 固定方案保留开发48组与另48组独立设计；新建设机制在原开发48组内配对比较\\
碳核算 & 年度历史有效映射加简化工业与固定碳汇；小时用IPCC燃烧系数 & 年度剩余净排衡量所选技术组合下的中和缺口；小时报告实际燃料燃烧CO$_2$\\
\bottomrule\end{tabularx}\end{table}\FloatBarrier
